% report/report.tex
\documentclass[11pt,a4paper]{article}

% --------- Packages ---------
\usepackage[margin=2.5cm]{geometry}
\usepackage[T1]{fontenc}
\usepackage[utf8]{inputenc}
\usepackage{lmodern}
\usepackage{microtype}
\usepackage{amsmath}
\usepackage{hyperref}
\usepackage{graphicx}
\usepackage{booktabs}
\usepackage{enumitem}
\usepackage{xcolor}
\usepackage{listings}

\hypersetup{
  colorlinks=true,
  linkcolor=blue,
  urlcolor=blue,
  citecolor=blue
}

% --------- Listings (AnB / OFMC outputs) ---------
\lstdefinestyle{anb}{
  basicstyle=\ttfamily\small,
  breaklines=true,
  frame=single,
  columns=fullflexible,
  showstringspaces=false
}

\lstdefinestyle{terminal}{
  basicstyle=\ttfamily\small,
  breaklines=true,
  frame=single,
  columns=fullflexible,
  showstringspaces=false
}

% --------- Simple helper command ---------
\newcommand{\sectionauthor}[1]{\noindent\textbf{Section author: #1}\par\medskip}

% --------- Title info ---------
\title{Open Authorization: Protocol Design and Formal Analysis with OFMC}
\author{
  Kristófer Birgir Hjörleifsson - s253155 \and Mikael Máni Eyfeld Clarke - s253024
}
\date{Hand-in: Mar 16, 2026}

\begin{document}
\maketitle

\begin{abstract}
We design and incrementally refine an open authorization protocol in which a user A delegates
read-only access to a photo resource at server B to a photo book service P, using an honest identity
provider IdP. The protocol is specified in AnB and analyzed using OFMC under the Dolev--Yao attacker
model. This report documents the development process, discovered attacks, design iterations, and the
final protocol, including weekly special tasks.
\end{abstract}

\tableofcontents
\newpage

% ------------------------------------------------
\section{Group Setup and Section Ownership}
\sectionauthor{Kristófer Birgir Hjörleifsson}

\begin{itemize}[leftmargin=*]
  \item \textbf{Group members:} Kristófer Birgir Hjörleifsson, Mikael Máni Eyfeld Clarke
  \item \textbf{Repository:} (https://github.com/kristoferbirgir/Project1-LogicForSecurity)
  \item \textbf{Submitted artifacts:} This report (PDF), AnB files, and the lab logbook.
\end{itemize}

\noindent\textbf{Section ownership requirement (DTU):} Each section in this report lists exactly one
author. Any section without a single designated author may be counted as not submitted.

% ------------------------------------------------
\section{Resources Used}
\sectionauthor{Kristófer Birgir Hjörleifsson}

\begin{itemize}[leftmargin=*]
  \item Course slides and lecture notes (weeks 2--7).
  \item OFMC / AnB documentation and examples from the course material.
  \item AVISPA project documentation for AnB syntax reference.
  \item AI tools: GitHub Copilot used for protocol development iteration and documentation.
  \item Group discussions between team members.
\end{itemize}

% ------------------------------------------------
\section{Scenario and Security Goals}
\sectionauthor{Mikael Máni Eyfeld Clarke}

\subsection{Scenario}
A stores holiday photos at server B. A wants to create a photo book using service P and wishes to
authorize P to access a subset of A's photos at B without downloading and re-uploading them.
A authenticates using an identity provider IdP via a shared password. The password is considered
guessable and must not be used as an encryption key. All parties except A have public/private keys.
All parties know IdP's public key, and IdP knows all public keys. IdP is assumed honest and trusted.

\subsection{Agents}
\begin{itemize}[leftmargin=*]
  \item \textbf{A (User):} Owns photos, authenticates via password, no cryptographic keys.
  \item \textbf{B (Photo Server):} Stores A's photos, has key pair \texttt{pk(B)/inv(pk(B))}.
  \item \textbf{P (Printing Service):} Wants to access A's photos, has key pair \texttt{pk(P)/inv(pk(P))}.
  \item \textbf{IdP (Identity Provider):} Authenticates A, issues authorization tokens, has key pair \texttt{pk(IdP)/inv(pk(IdP))}.
\end{itemize}

\subsection{Security and Functional Goals}
We target the following goals:
\begin{itemize}[leftmargin=*]
  \item \textbf{Authorization:} B should only release photos if authorized by A via IdP.
  \item \textbf{Password secrecy:} A's password must not be revealed to P, B, or the intruder.
  \item \textbf{Photo confidentiality:} Photos should remain secret between authorized parties.
  \item \textbf{Authentication:} B must be able to verify IdP signed the authorization token.
\end{itemize}

% ------------------------------------------------
\section{Week 2: Informal Design and Static Dolev--Yao Analysis}
\sectionauthor{Kristófer Birgir Hjörleifsson}

\subsection{Informal Outline}
The basic protocol flow:
\begin{enumerate}
  \item A requests authorization from IdP (encrypted with IdP's public key)
  \item IdP creates signed token authorizing P to access A's photos at B
  \item A forwards token to P
  \item P presents token to B (encrypted with B's public key)
  \item B verifies token and sends photos to P
\end{enumerate}

\subsection{First AnB Model (Simplified)}
Initial simplification: A has a key pair (temporary, removed in Week 3).

\begin{lstlisting}[style=anb,caption={Week 2 protocol excerpt}]
Actions:
  A -> IdP: {A, P, B}(pk(IdP))
  IdP -> A: {A, P, B}(inv(pk(IdP)))
  A -> P: {A, P, B}(inv(pk(IdP)))
  P -> B: {{A, P, B}(inv(pk(IdP)))}(pk(B))
  B -> P: {photos(A)}(pk(P))
\end{lstlisting}

\subsection{OFMC Result}
OFMC found a \textbf{signature forgery attack}: the intruder creates a token signed with their own key, and B cannot distinguish \texttt{pk(IdP)} from \texttt{pk(i)}.

\subsection{Special Task: Static Analysis (Passive Intruder)}
\textbf{Assumption:} Intruder only eavesdrops (passive), does not inject messages.

\textbf{What intruder observes:}
\begin{itemize}[leftmargin=*]
  \item Message 1: \texttt{\{A, P, B\}(pk(IdP))} -- cannot decrypt (no \texttt{inv(pk(IdP))})
  \item Message 2: \texttt{\{A, P, B\}(inv(pk(IdP)))} -- can verify signature, learns A, P, B
  \item Message 3: Same signed token -- learns nothing new
  \item Message 4: \texttt{\{\{token\}\}(pk(B))} -- cannot decrypt outer layer
  \item Message 5: \texttt{\{photos(A)\}(pk(P))} -- cannot decrypt
\end{itemize}

\textbf{Intruder's knowledge (passive):} Agent identities A, P, B, IdP (from signed token). \textbf{Cannot learn:} photos, password, or private keys.

\textbf{Conclusion:} Protocol is secure against passive eavesdropping. The signature forgery attack requires active message injection.

% ------------------------------------------------
\section{Week 3: Identity Provider and Lazy Intruder Executability}
\sectionauthor{Mikael Máni Eyfeld Clarke}

\subsection{Protocol Refinement}
\textbf{Key change:} A no longer has a key pair. A authenticates to IdP using password.

\begin{lstlisting}[style=anb,caption={Week 3: Password authentication}]
A -> IdP: {A, P, B, pwd(A)}(pk(IdP))
\end{lstlisting}

The password is encrypted under IdP's public key, so the intruder cannot learn it. Password is NOT used as an encryption key -- only as a credential.

\subsection{Special Task: Lazy Intruder Method (P Executable by Intruder)}
We show role P is executable by constructing all outgoing messages from intruder's knowledge.

\textbf{Role P receives and sends:}
\begin{enumerate}
  \item Receives: \texttt{token = \{IdP, A, P, B\}(inv(pk(IdP)))}
  \item Sends: \texttt{\{P, token\}(pk(B))}
\end{enumerate}

\textbf{Lazy Intruder Analysis:}
\begin{itemize}[leftmargin=*]
  \item \textbf{Step 1:} Intruder intercepts message 3 (A $\to$ P), learns \texttt{token}
  \item \textbf{Step 2:} Intruder knows \texttt{P} (public agent name)
  \item \textbf{Step 3:} Intruder knows \texttt{pk(B)} (public key, in initial knowledge)
  \item \textbf{Step 4:} Intruder computes \texttt{\{P, token\}(pk(B))} using encryption rule
\end{itemize}

\textbf{Conclusion:} Role P is executable by the intruder. Every message P sends can be constructed from publicly available information and intercepted messages. This is expected -- P acts as a relay.

% ------------------------------------------------
\section{Week 4: Typing, Formats, and Public-Key Lookup}
\sectionauthor{Kristófer Birgir Hjörleifsson}

\subsection{Type-Flaw Resistance via Formats}
Type flaws occur when an intruder substitutes a message of one type where another is expected. We prevent this by:
\begin{itemize}[leftmargin=*]
  \item Including agent identities in messages as implicit format tags
  \item Using different tuple structures for different message types
\end{itemize}

\begin{lstlisting}[style=anb,caption={Message formats}]
Msg 1: {A, P, B, pwd(A)}(pk(IdP))       % 4-tuple with password
Msg 2: {IdP, A, P, B}(inv(pk(IdP)))     % 4-tuple, agents only
Msg 4: {P, {token}(inv(pk(IdP)))}(pk(B)) % Nested structure
Msg 5: {B, photos(A)}(pk(P))            % 2-tuple
\end{lstlisting}

\subsection{Public-Key Lookup Subprotocol}
A separate protocol allows A to obtain public keys from IdP:

\begin{lstlisting}[style=anb,caption={Key lookup protocol}]
A -> IdP: {A, B, pwd(A)}(pk(IdP))
IdP -> A: {B, pk(B)}(inv(pk(IdP)))
\end{lstlisting}

\subsection{Special Task: Type-Flaw Resistance Proof}
\textbf{Definition:} Protocol is type-flaw resistant if for any two non-variable messages $M_1$, $M_2$ in SMP, if they have a unifier $\sigma$, then $\text{type}(M_1) = \text{type}(M_2)$.

\textbf{Analysis of message types:}
\begin{center}
\begin{tabular}{lll}
\toprule
Message & Structure & Type \\
\midrule
Msg 1 & (Agent, Agent, Agent, Function) & AuthRequest \\
Msg 2 & (Agent, Agent, Agent, Agent) & Token \\
Msg 4 outer & (Agent, SignedMsg) & PhotoRequest \\
Msg 5 & (Agent, Function) & PhotoResponse \\
\bottomrule
\end{tabular}
\end{center}

\textbf{Unification check:}
\begin{itemize}[leftmargin=*]
  \item Msg 1 vs Msg 2: Different types (\texttt{pwd(A)} $\neq$ Agent) -- no unifier
  \item Msg 4 vs Msg 5: Different types (SignedMsg $\neq$ Function) -- no unifier
\end{itemize}

\textbf{Conclusion:} Protocol is type-flaw resistant.

% ------------------------------------------------
\section{Week 5: Channels (TLS-like Pseudonymous Channels)}
\sectionauthor{Mikael Máni Eyfeld Clarke}

\subsection{Channel Notation}
\begin{center}
\begin{tabular}{ll}
\toprule
Notation & Meaning \\
\midrule
\texttt{A -> B: M} & Insecure (Dolev-Yao) \\
\texttt{A *-> B: M} & Confidential only (TLS, client not auth) \\
\texttt{A ->* B: M} & Authenticated only (signed) \\
\texttt{A *->* B: M} & Secure channel (mutual TLS) \\
\bottomrule
\end{tabular}
\end{center}

\subsection{Channel Variant}
\begin{lstlisting}[style=anb,caption={Week 5: TLS channels}]
A *->* IdP: A, P, B, pwd(A)           % TLS to IdP
IdP *->* A: {IdP, A, P, B}(inv(pk(IdP)))  % Signed token
A *-> P: {token}(inv(pk(IdP)))        % Forward token
P *->* B: P, {token}(inv(pk(IdP)))    % Present to B
B *->* P: B, photos(A)                % Response
\end{lstlisting}

\subsection{Special Task: What Was Replaced}
\begin{center}
\begin{tabular}{lll}
\toprule
Original & Channel Version & Replaced? \\
\midrule
\texttt{\{M\}(pk(IdP))} & \texttt{A *-> IdP: M} & Yes \\
\texttt{\{M\}(pk(B))} & \texttt{P *-> B: M} & Yes \\
\texttt{\{M\}(pk(P))} & \texttt{B *-> P: M} & Yes \\
\texttt{\{M\}(inv(pk(IdP)))} & Cannot replace & \textbf{No} \\
\bottomrule
\end{tabular}
\end{center}

\textbf{Why signature cannot be replaced:} B needs to verify IdP approved the authorization, but B has no direct TLS channel to IdP. The token travels through A and P, so B must verify the signature cryptographically.

% ------------------------------------------------
\section{Week 6: Privacy and Guessable Password}
\sectionauthor{Kristófer Birgir Hjörleifsson}

\subsection{Guessable Secrets}
A guessable secret (like a password) can be subject to offline dictionary attacks. The intruder captures encrypted data containing the password and tries decrypting with guessed passwords.

\subsection{Security with Guessable Password}
Our protocol protects against guessing attacks because:
\begin{itemize}[leftmargin=*]
  \item Password is sent over secure channel (\texttt{A *->* IdP})
  \item Intruder never sees ciphertext containing password
  \item No offline verification possible
\end{itemize}

\textbf{OFMC Goal:} \texttt{pwd(A) guessable secret between A, IdP}

\textbf{Result:} No guessing attack found (only signature forgery under dishonest IdP).

\subsection{Special Task: Insecure Variant Demonstration}
We created an insecure variant to demonstrate the attack:

\begin{lstlisting}[style=anb,caption={Insecure version (vulnerable to guessing)}]
A -> IdP: {A, P, B, pwd(A)}(pk(IdP))   % Regular encryption
\end{lstlisting}

\textbf{OFMC Result:} \texttt{ATTACK\_FOUND (guesswhat)}
\begin{lstlisting}[style=terminal,caption={Guessing attack trace}]
(A,1) -> i: {A,P,B,pwd(A)}_(pk(IdP))
i can produce secret {A,P,B,guessPW}_(pk(IdP))
\end{lstlisting}

\textbf{Conclusion:} Encryption alone does not protect guessable secrets. Secure channels (TLS) prevent offline guessing by hiding the ciphertext entirely.

% ------------------------------------------------
\section{Final Protocol (AnB Specification Summary)}
\sectionauthor{Mikael Máni Eyfeld Clarke}

\subsection{Final Message Flow}
\begin{lstlisting}[style=anb,caption={Final protocol (photo\_auth\_final.AnB)}]
A *->* IdP: A, P, B, pwd(A)
IdP *->* A: {IdP, A, P, B}(inv(pk(IdP)))
A *-> P: {IdP, A, P, B}(inv(pk(IdP)))
P *->* B: P, {IdP, A, P, B}(inv(pk(IdP)))
B *->* P: B, photos(A)
\end{lstlisting}

\textbf{Step-by-step:}
\begin{enumerate}
  \item A authenticates to IdP with password over TLS
  \item IdP creates signed authorization token
  \item A forwards token to P over confidential channel
  \item P presents token to B with mutual TLS
  \item B verifies token signature and sends photos to P
\end{enumerate}

\subsection{Initial Knowledge}
\begin{center}
\begin{tabular}{ll}
\toprule
Agent & Knowledge \\
\midrule
A & A, B, P, IdP, pwd(A) \\
B & A, B, P, IdP, pk(B), pk(IdP), inv(pk(B)), photos(A) \\
P & A, B, P, IdP, pk(B), pk(P), pk(IdP), inv(pk(P)) \\
IdP & A, B, P, IdP, pk(B), pk(P), pk(IdP), inv(pk(IdP)), pwd(A) \\
\bottomrule
\end{tabular}
\end{center}

\subsection{Goals}
\begin{itemize}[leftmargin=*]
  \item \texttt{B authenticates IdP on IdP, A, P, B} -- B verifies IdP signed token
  \item \texttt{photos(A) secret between B, P} -- Photos only known to B and P
  \item \texttt{pwd(A) guessable secret between A, IdP} -- Password protected
\end{itemize}

\subsection{Verification Results}
\textbf{OFMC Result:} ATTACK\_FOUND (signature forgery)

The attack involves the intruder acting as IdP, signing a forged token with their own key. This attack is \textbf{excluded by assumption} since IdP is assumed honest.

\textbf{Under honest IdP assumption:} Protocol is secure.

% ------------------------------------------------
\section{Assumptions, Simplifications, and Limitations}
\sectionauthor{Kristófer Birgir Hjörleifsson}

\subsection{Assumptions}
\begin{itemize}[leftmargin=*]
  \item \textbf{IdP is honest:} The identity provider never acts maliciously or is compromised. Attacks involving dishonest IdP are excluded.
  \item \textbf{Secure channels abstract TLS:} We assume TLS provides confidentiality and authentication as modeled.
  \item \textbf{Keys are pre-distributed:} Servers know each other's public keys.
\end{itemize}

\subsection{Simplifications}
\begin{itemize}[leftmargin=*]
  \item \textbf{Photos as abstract secret:} Modeled as \texttt{photos(A)} rather than actual data.
  \item \textbf{No scope/folder:} Token authorizes access to all photos, not a specific subset.
  \item \textbf{No timestamps/nonces:} Token has no expiration or freshness mechanism.
  \item \textbf{Single session:} OFMC verifies bounded number of sessions.
\end{itemize}

\subsection{Limitations}
\begin{itemize}[leftmargin=*]
  \item \textbf{Signature forgery persists:} B cannot distinguish pk(IdP) from pk(i) without out-of-band trust establishment.
  \item \textbf{No revocation:} Once issued, tokens cannot be revoked.
  \item \textbf{IdP single point of failure:} If IdP is compromised, all security fails.
\end{itemize}

% ------------------------------------------------
\section{Conclusion}
\sectionauthor{Mikael Máni Eyfeld Clarke}

We designed an open authorization protocol allowing user A to authorize printing service P to access photos at server B, using identity provider IdP for authentication.

\textbf{Main attack encountered:} Signature forgery -- the intruder creates tokens signed with their own key. This is a fundamental Dolev-Yao limitation where B cannot distinguish honest from dishonest signers without pre-established trust.

\textbf{Solution:} We assume IdP is honest (as per assignment). Under this assumption, the protocol provides:
\begin{itemize}[leftmargin=*]
  \item Authentication of IdP's authorization via signed tokens
  \item Confidentiality of photos (only B and P)
  \item Protection of guessable password via secure channels
\end{itemize}

\textbf{Final protocol uses:} TLS channels for confidentiality, cryptographic signatures for authorization tokens that travel through untrusted intermediaries.

% ------------------------------------------------
\appendix
\section{Submission Map (Where to Find Things)}
\sectionauthor{Kristófer Birgir Hjörleifsson}

\begin{itemize}[leftmargin=*]
  \item Lab logbook: \texttt{logbook/logbook.md}
  \item AnB protocol versions:
    \begin{itemize}
      \item \texttt{anb/week2\_v1.AnB} -- Basic protocol
      \item \texttt{anb/week3\_v1.AnB} -- Password authentication
      \item \texttt{anb/week4\_v1.AnB} -- Type-flaw resistant
      \item \texttt{anb/week5\_v1.AnB} -- TLS channels
      \item \texttt{anb/week5\_v1\_tls.AnB} -- Server-side only TLS
      \item \texttt{anb/week6\_v1.AnB} -- Guessable password
      \item \texttt{anb/week6\_insecure.AnB} -- Insecure variant (demo)
    \end{itemize}
  \item Final AnB files:
    \begin{itemize}
      \item \texttt{anb/final/photo\_auth\_final.AnB} -- Final protocol
      \item \texttt{anb/final/key\_lookup.AnB} -- Key lookup subprotocol
    \end{itemize}
\end{itemize}

\end{document}